% Part: counterfactuals
% Chapter: minimal-change-semantics
% Section: transitivity

\documentclass[../../../include/open-logic-section]{subfiles}

\begin{document}

\olfileid{cnt}{min}{cpo}

\olsection{换质位律}

实质条件句和严格条件句都等价于它们的换质位式。反事实条件句则不然。以下是 Kratzer（克拉策尔）给出的一个例子：
\begin{quote}
  如果歌德没有在1832年去世，那么他现在（依然）是死的。

  如果歌德现在没死，那么他就会是在1832年去世的。
\end{quote}
第一句话是真的：人活不到几百年。第二句话显然是假的：如果歌德现在没死，那他就还活着，因此不可能在1832年去世。

\begin{ex}\ollabel{ex:contraposition-counterex}
  球系语义使换质位律无效，即我们有 $p \cif q \Entails/ \lnot q \cif \lnot p$。将 $p$ 理解为“歌德没有在1832年去世”，将 $q$ 理解为“歌德现在死了”。我们可以在一个模型 $\mModel{M_1} = \tuple{W, O, V}$ 中刻画这一点，其中 $W =
  \{w, w_1, w_2\}$，$O = \{\{w\}, \{w, w_1\}, \{w, w_1, w_2\}\}$，$V(p) = \{w_1, w_2\}$，且 $V(q) = \{w, w_1\}$。于是，$w$ 是现实世界，歌德在1832年去世且现在仍是死的；$w_1$ 是一个（较近的）世界，歌德比如在1833年去世，且现在仍是死的；而 $w_2$ 是一个（较远的）世界，歌德仍然活着。
\begin{figure}
\begin{center}
\begin{tikzpicture}[scale=.6,layerwidth=1.5]\tiny
  \spheresystem{3}
  \begin{scope}
    \clip \propositionplot{0}{.8}{50}{4.5} -- (4.5,-4.5) -- (-4.5,-4.5) -- (-4.5,4.5) -- (4.5,4.5);
    \spherelayer{3}
  \end{scope}
  \propositionintersect{0}{2}{110}{5.5}
  \proposition{0}{.8}{50}{4.5}
  \path (0,0) node[world] {$w$};
  \spherepos{0}{2}{node[world] {$w_1$}}
  \spherepos{-50}{3}{node[world] {$w_2$}}
  \spherepos{28}{3}{node {$q$}}
  \spherepos{42}{3}{node {$\lnot q$}}
  \spherepos{-55}{4}{node {$p$}}
  \spherepos{-67}{4}{node {$\lnot p$}}
\end{tikzpicture}
\caption{换质位的反例}
\ollabel{fig:contraposition}
\end{center}
\end{figure}
存在一个容纳 $p$ 的球层 $S =
\{w, w_1\}$，且 $p \lif q$ 在其中的所有世界上为真，因此 $\mSat{M}{p
  \cif q}[w]$。然而，容纳 $\lnot q$ 的球层 $\{w, w_1,
w_2\}$ 包含一个世界，即~$w_2$，在该世界中 $q$ 为假且 $p$ 为真，所以 $\mSat/{M}{\lnot q \lif \lnot p}[w_2]$。
\end{ex}

\end{document}